\documentclass[12pt]{amsart}
%   \overfullrule=5pt
%\usepackage{fullpage}
\usepackage{etex} % Patrick: I loaded this package because otherwise I get obscure error messages
\usepackage[active]{srcltx}
\usepackage{calc,amssymb,amsthm,amsmath,amscd, eucal,ulem}
\usepackage{alltt}
%\usepackage[left=1in,top=1in,right=1in,bottom=1in]{geometry}
\synctex=1
%\usepackage{mathtools}
\RequirePackage[dvipsnames,usenames]{color}
\let\mffont=\sf
\normalem
\input{mabliautoref.sty}
\input{xy}
\xyoption{all}
\input{kmacros3.sty}
\usepackage{tikz}
\usepackage{amsfonts, mathrsfs}
\usepackage{calligra}
\usepackage{stmaryrd} %power series bracket
\usepackage{dcpic, pictexwd}
%\usepackage{pdfsync}
\usepackage[left=1in,top=1in,right=1in,bottom=1in]{geometry}

\numberwithin{equation}{theorem}

%\renewcommand{\thefootnote}{\fnsymbol{footnote}}
\newcommand{\D}{\displaystyle}
\newcommand{\til}{\widetilde}
\newcommand{\ol}{\overline}
\newcommand{\F}{\mathbb{F}}
\DeclareMathOperator{\E}{E}
\renewcommand{\:}{\colon}
\newcommand{\eg}{{\itshape e.g.} }
\renewcommand{\m}{\mathfrak{m}}
\renewcommand{\n}{\mathfrak{n}}
\newcommand{\Tt}{{\mathfrak{T}}}
\newcommand{\calC}{\mathcal{C}}
\newcommand{\RamiT}{\mathcal{R_{T}}}
\DeclareMathOperator{\edim}{edim}
\DeclareMathOperator{\length}{length}
\DeclareMathOperator{\monomials}{monomials}
\DeclareMathOperator{\Fitt}{Fitt}
\DeclareMathOperator{\depth}{depth}
\newcommand{\Frob}[2]{{#1}^{1/p^{#2}}}
\newcommand{\Frobp}[2]{{(#1)}^{1/p^{#2}}}
\newcommand{\FrobP}[2]{{\left(#1\right)}^{1/p^{#2}}}
\newcommand{\extends}{extends over $\Tt${}}
\newcommand{\extension}{extension over $\Tt${}}
\newcommand{\piet}{\pi_1^{\textnormal{\'{e}t}}} % \'Etale fundamental group
\newcommand{\FEt}{\textsf{F\'{E}t}} % Category of finite \'etale covers
\DeclareMathOperator{\ns}{ns}
%\newcommand{\Spec}{\text{Spec}} % Spectrum
\newcommand{\pSpec}{\text{Spec}^{\circ}} % Puntured Spectrum
\newcommand{\etInCdOne}{\etale{} in codimension $1$}
\DeclareMathOperator{\DIV}{Div}

\usepackage{setspace}
\usepackage{hyperref}
\usepackage{enumerate}
\usepackage{graphicx}
\usepackage[all,cmtip]{xy}

\usepackage{verbatim}

\theoremstyle{theorem}
\newtheorem*{mainthm}{Main Theorem}
\newtheorem*{mainthma}{Main Theorem A}
\newtheorem*{mainthmb}{Main Theorem B}
\newtheorem*{mainthmc}{Main Theorem C}
\newtheorem*{atheorem}{Theorem}

% The todo box!
\def\todo#1{\textcolor{Mahogany}%
{\footnotesize\newline{\color{Mahogany}\fbox{\parbox{\textwidth-15pt}{\textbf{todo:
} #1}}}\newline}}

% What is going on?  Why isn't \O a script O?!?
\renewcommand{\O}{\mathscr O}

\begin{document}
\title{\'Etale fundamental groups of strongly $F$-regular schemes}
\author[B.~Bhatt]{Bhargav Bhatt}
\author[J.~Carvajal-Rojas]{Javier Carvajal-Rojas}
\author[P.~Graf]{Patrick Graf}
\author[K.~Schwede]{Karl Schwede}
\author[K.~Tucker]{Kevin Tucker}
\address{Department of Mathematics\\ University of Michigan\\Ann Arbor\\ MI 48109}
\email{\href{mailto:bhargav.bhatt@gmail.com}{bhargav.bhatt@gmail.com}}
\address{Department of Mathematics\\ University of Utah\\ Salt Lake City\\ UT 84112 \newline\indent
Escuela de Matem\'atica\\ Universidad de Costa Rica\\ San Jos\'e 11501\\ Costa Rica}
\email{\href{mailto:carvajal@math.utah.edu}{carvajal@math.utah.edu}}
\address{Lehrstuhl f\"ur Mathematik I, Universit\"at Bay\-reuth, 95440 Bayreuth, Germany}
\email{\href{mailto:patrick.graf@uni-bayreuth.de}{patrick.graf@uni-bayreuth.de}}
\address{Department of Mathematics\\ University of Utah\\ Salt Lake City\\ UT 84112}
\email{\href{mailto:schwede@math.utah.edu}{schwede@math.utah.edu}}
\address{Department of Mathematics\\ University of Illinois at Chicago\\Chicago\\ IL 60607}
\email{\href{mailto:kftucker@uic.edu}{kftucker@uic.edu}}

\thanks{Bhatt was supported in part by the NSF Grant DMS \#1501461 and a Packard fellowship}
\thanks{Carvajal-Rojas was supported in part by the NSF FRG Grant DMS \#1265261/1501115}
\thanks{Graf was supported in part by the DFG grant ``Zur Positivit\"at in der komplexen
Geometrie''}
\thanks{Schwede was supported in part by the NSF FRG Grant DMS \#1265261/1501115 and NSF CAREER Grant DMS \#1252860/1501102}
\thanks{Tucker was supported in part by NSF Grant DMS \#1602070 and a fellowship from the Sloan foundation.}

\subjclass[2010]{14F35, 14F20, 14E20, 14F18, 14B05, 13A35, 13B40}

%\date{\today}

\begin{abstract}
We prove that a strongly $F$-regular scheme $X$ admits a finite, generically Galois, and \'etale-in-codimension-one  cover $\tld X \to X$ such that the \'etale fundamental groups of $\tld X$ and $\tld X_{\reg}$ agree.
Equivalently, every finite \'etale cover of $\tld X_{\reg}$ extends to a finite \'etale cover of $\tld X$.
This is analogous to a result for complex klt varieties by Greb, Kebekus and Peternell.
\end{abstract}

\maketitle

\section{Introduction}

In \cite{XuFinitenessOfFundGroups}, Chenyang Xu proved that the algebraic local fundamental group (the profinite completion of the topological fundamental group of the link) of a complex klt singularity is finite.
This result was then used by Greb--Kebekus--Peternell \cite{GrebKebekusPeternellEtaleFundamental} to show that for a complex quasi-projective variety $X$ with klt singularities, any sequence of finite quasi-\'etale (\ie~\'etale in codimension one) and generically Galois covers eventually becomes \etale.
In particular, for any quasi-projective klt variety $X/\C$ there exists a finite quasi-\'etale generically Galois cover $\rho\colon \tld X \to X$ such that any further quasi-\'etale cover $Y \to \tld X$ is actually \'etale.

Inspired by the relation between klt singularities and strongly $F$-regular singularities in characteristic $p > 0$ \cite{HaraWatanabeFRegFPure}, a subset of the authors showed that if $(R, \fram)$ is a strictly Henselian strongly $F$-regular local ring and $U \subseteq \Spec R$ is the regular locus, then $\pi_1^{\et}(U)$ is finite \cite{CarvajalSchwedeTuckerEtaleFundFsignature}, a variant of Xu's result in characteristic $p > 0$.  In this paper we use this result, a theorem of Gabber \cite{TravauxdeGabber}, and recent work on the non-local behavior of $F$-signature \cite{deStefaniPolstraYao-Globalizing}, to prove a variant of the result of Greb--Kebekus--Peternell in characteristic $p > 0$.

Philosophically, both these results are studying the obstructions of extending finite \'etale covers of the regular locus $X_{\reg}$ of a scheme $X$ to the whole scheme.  The expected general answer is that for more severe singularities on $X$, there are more obstructions.
Analogous to \cite{GrebKebekusPeternellEtaleFundamental}, we show that schemes with strongly $F$-regular singularities are mild in this sense.

\begin{mainthm*}[\textnormal{\autoref{thm.QuasiEtaleStabilize}, \autoref{cor.MaximalQuasiEtale}}]
Suppose $X$ is an $F$-finite Noetherian integral strongly $F$-regular scheme.
Suppose we are given a sequence of finite surjective quasi-\'etale (\'etale in codimension one) morphisms of normal integral schemes
\[ X = X_0 \xleftarrow{\;\;\gamma_1\;\;} X_1 \xleftarrow{\;\;\gamma_2\;\;} X_2 \xleftarrow{\;\;\gamma_3\;\;} \cdots \]
such that each $X_i/X$ is generically Galois.
Then all but finitely many of the $\gamma_i$ are \'etale.

In particular, there exists a finite quasi-\'etale generically Galois cover $\tld X \to X$ so that any further finite quasi-\'etale cover $Y \to \tld X$ is actually \'etale.
Equivalently, the map $\piet(\tld X_{\reg}) \to \piet(\tld X)$ induced by the inclusion of the regular locus is an isomorphism.
\end{mainthm*}

Note that we do not require any quasi-projectivity hypothesis.
We actually obtain a slightly stronger version, just as in \cite{GrebKebekusPeternellEtaleFundamental}.
As a corollary, we also obtain a variant of \cite[Theorem 1.10]{GrebKebekusPeternellEtaleFundamental}, see \autoref{cor.SimultaneousIndex1Cover}, which says in particular that there exists an integer $N > 0$ such that for every $\bQ$-Cartier divisor $D$ on $X$ with index not divisible by $p$, $N \cdot D$ is in fact Cartier.
Finally, we observe that the map $\tld X \to X$ from our Main~Theorem above is tame wherever it is not \'etale, see \autoref{cor.MaximalTameCover}.

\subsection*{Acknowledgements}

The authors began working on this project while visiting CIRM Luminy in September 2016.

\section{Notation and conventions}

\begin{convention}
All rings and schemes considered will be Noetherian.  We will typically also be working with normal connected and hence integral schemes.
%Thus when considering \'etale fundamental groups, we will always choose the base point of a scheme $X$ to be a fixed geometric \emph{generic} point, that is, $\Spec K(X)^{\text{sep}} \to X$ for a fixed separable closure of the fraction field $K(X)$.
\end{convention}

%\begin{definition}
%Suppose $X$ is a Noetherian scheme and $W \to X$ is an affine $X$-scheme.  We define $\O_{X,W}^{\sh}$, the \emph{coordinate ring of the strict henselization of $X$ along $W$}, to be the direct limit $\varinjlim \Gamma(Y, \O_Y)$ of the coordinate rings of $Y$ over all \'etale maps $Y\to X$ which factor $W \to X$ (that is, $W \xrightarrow{\;\exists\;} Y \to X$).
%\end{definition}

%Note that there is a natural map $\Spec \O_{X, W}^{\sh} \to X$ which factors $W \to X$.

\begin{notation}
Let $\rho\colon Y \to X$ be a morphism of normal schemes.
\begin{enumerate}
\item
We say that $\rho$ is \emph{quasi-\'etale} if it is quasi-finite and there exists a closed subset $Z \subseteq Y$ of codimension $\ge 2$ such that $\rho|_{Y \setminus Z}\colon Y \setminus Z \to X$ is \'etale.
Notice that if $\rho$ is further finite, then $Y$ coincides with the integral closure of $X$ in $K(Y)$.
\item
\label{not.b.branchlocus}
Assume that $\rho$ is finite and surjective. Let $W \subseteq X$ be the largest Zariski-open subset such that $\rho^{-1}(W) \to W$ is \'etale.
The \emph{branch locus} of $\rho$ is defined to be $X \setminus W$.
It is worth to remark that by Purity of the Branch Locus, if $\rho$ is quasi-\'etale then its branch locus is contained in the non-regular locus of $X$.
\item
If $X$ and $Y$ are connected, we say that $\rho$ finite and surjective is \emph{generically Galois} if the fraction fields extension $K(Y)/K(X)$ is Galois.
In this case, it follows that $\rho^{-1}(W) \to W$ is Galois, where $W$ is as in \autoref{not.b.branchlocus}.
\item
By a \emph{geometric point} $\bar x \in X$ of a scheme $X$ we mean a morphism $\Spec L \to X$, where $L$ is a separably closed field.
Note that $\bar x \to X$ factors through $\Spec k(x) \to X$, where $x \in X$ is the scheme-theoretic image point.
\item
Suppose that $\bar s, \bar t \in X$ are geometric points of a scheme $X$.
A \emph{specialization of geometric points $\bar s \rightsquigarrow \bar t$} is a factorization
\[
\bar s \to \Spec \O_{X, \overline{t}}^{\sh} \to X
\]
where the composition is the map defining the geometric point $\bar s \to X$.
Note that then we get a map $\Spec \O_{X,\overline{s}}^{\sh} \to \Spec \O_{X,\overline{t}}^{\sh}$.
In particular, the image of $\overline{t}$ as a scheme theoretic point on $X$ is contained in the Zariski closure of the image of $\overline{s}$.
Conversely, if $\bar s \in X$ is a geometric point and $t \in \overline{ \{ s \} }$, then we can find a geometric point $\bar t \in X$ with scheme-theoretic image point $t$ such that $\bar s$ specializes to $\bar t$.
\end{enumerate}
\end{notation}

\begin{comment}
{\color{blue} Javier: This is what I understand by a specialization of geometric points, I hope it coincides with what Bharghav said when we were skyping the other day. I also apologize for my so wordy definitions (and poor grammar), feel free to cut all you want.}
By a \emph{geometric point} $\bar x \to X$ of a scheme $X$ we mean a morphism $\Spec k_x \to X$ that factors through $\Spec k(x) \to X$, where $x\in X$ is a scheme-theoretic point and $k(x) \subseteq k_x$ is a separably closed extension of fields. In the rest of this article we often forget the ``bar'' in the notation and deliberately identify the geometric points with its image at $X$. Given two geometric points $\bar{s}\to X$, $\bar{t} \to X$  of $X$, we say $\bar{s}$ \emph{specializes to} $\bar{t}$ (alternatively, $\bar{t}$ \emph{is a specialization of} $\bar{s}$) and denote it by $\bar{s} \rightsquigarrow \bar{t}$, if there exists $V$ a valuation ring of $k_s$ together with a morphism $\Spec V \to X$ such that there exists a commutative diagram:
\[
\xymatrix{
\bar{t} \ar[rd]_-{\exists} \ar[rrd]\\
& \Spec V \ar[r] & X\\
\bar{s} \ar[ru] \ar[rru]
}
\]
\end{enumerate}
It is worth noting that if $\bar{s} \rightsquigarrow \bar{t}$ then $t$ belongs to the Zariski-closure of $\{s\}$, say $\overline{\{s\}}$. The converse is not true, nonetheless it holds by suitably choosing $k_t$ (for given $k_s$). More precisely, given a geometric point $\bar{s}$ and $t \in \overline{\{s\}}$, there is a separably closed extension of $k(t)$ such that the associated geometric point specializes $\bar{s}$.  Indeed, if we denote by $Z$ the Zariski-closure of $\{s\}\subseteq X$ with its reduced induced scheme structure, we have $\O_{Z, t} \subseteq K(Z)=k(s) \subseteq k_s$, in such a way that there is a valuation ring $V$ of $k_s$ dominating\footnote{Meaning a local inclusion $\O_{Z,t} \subseteq V$ of local rings} $\O_{Z,t}$, with inclusion of residue fields $k(t)\subseteq k(V)$. The composite
\[\Spec V \to \Spec \O_{Z,t} \to Z \to X\]
then gives the desired specialization, as the residue field $k(V)$ is separably closed {\color{blue} Javier: I'm not sure on how to argue this anymore, I'd like to use \cite[Tag~0BSQ]{stacks-project} as below, but I'm not seeing why $V$ is or can be taken normal. Patrick: I believe valuation rings are always normal? Atiyah--MacDonald, Prop.~5.18? \\
So you are taking $k_t = k(V)$, right? And then you want to argue that $V$ is strictly henselian, so we get the maps as in \autoref{prop.GabberStratification}?}

\todo{{\bf Karl:} I thought we wanted something much much simpler.  (Also, I'm not sure if we should define geometric point).  Anyways, I thought we just wanted the following.\vskip 9pt
If $s, t$ are geometric points, a \emph{specialization of geometric points} $s \rightsquigarrow t$ is a factorization
\[
s \to \Spec \O_{X, t}^{\sh} \to X.
\]
where the composite is the given geometric point map $s \to X$.
I think this is all we wanted, right?  I thought we wanted to avoid all valuations etc.}
\end{comment}

\section{Gabber's constructibility and applications to the fundamental group}

Recall the following theorem of Gabber.

\begin{theorem}\textnormal{\cite[Expos\'e XXI, Th\'eor\`eme 1.1, 1.3]{TravauxdeGabber}}
\label{thm.GabberConstructibility}
Suppose that $X$ is a normal excellent scheme and $Z \subseteq X$ a closed subscheme of codimension $\geq 2$.
Let $j\colon U = X \setminus Z \to X$ be the inclusion.
Then for all finite groups $G$, $R^i j_* (G_U)$ is constructible for $i = 0, 1$.
Here $G_U$ denotes the constant sheaf $G$ on $U$ equipped with the \'etale topology.
\end{theorem}

We use this to obtain:

\begin{proposition} \label{prop.GabberStratification}
Let $i\colon Z \hookrightarrow X$ be a closed subscheme of codimension $\ge 2$ in a normal excellent scheme.
Let $j\colon U = X \setminus Z \hookrightarrow X$ be the complement.
Suppose that
\[ \big| \pi_1^{\et}(\Spec \O_{X,x}^{\sh} \setminus {Z}^{\sh}) \big| \]
is bounded for all geometric points $x \in X$ (here $Z^{\sh} = Z_{X,x}^{\sh}$ is the inverse image of $Z$).
Then there exists a finite stratification $\{ Z_i \}_{i \in I}$ of $X$ into locally closed subschemes $Z_i$ with the following property:
\begin{itemize}
\item[] Fix a specialization $s \rightsquigarrow t$ of geometric points of $Z_i$. % and choose a valuation ring $V$ with algebraically %({\color{blue} separably})
%closed fraction field along with a map $\Spec V \to Z_i$ such that the closed point of $V$ maps to $t$ and the generic point maps to $s$.
Consider the map
    \begin{equation}
    \label{eq.MapOfHenselizations}
     \Spec(\O^{\sh}_{X,s}) \longrightarrow \Spec(\O^{\sh}_{X,t}).
    \end{equation}
    Remove the inverse images of $Z$ to obtain
    \begin{equation}
    \label{eq.TubularMapOfHenselizations}
    T_s \xrightarrow{\;\;\alpha\;\;} T_t.
    \end{equation}
    Then applying $\piet(-)$ to \autoref{eq.TubularMapOfHenselizations} yields an isomorphism.
\end{itemize}
\end{proposition}

\begin{proof}
Note there are only finitely many finite groups $G_x = \pi_1^{\et}(\Spec \O_{X,x}^{\sh} \setminus Z^{\sh})$ for geometric points $x \in X$, since there are only finitely many groups of size less than a given number.
Letting $G_x$ also denote the constant sheaf on $U_{\et}$, since $R^1 j_* G_x$ is constructible by \autoref{thm.GabberConstructibility}, we may choose a stratification $\{ Z_i \}$ of locally closed irreducible subschemes of $X$ so that $R^1 j_* G_x$ is locally constant on each $Z_i$ for all geometric points $x$. Note $R^1 j_* G_x$ is zero over $U$ so the induced stratification has the same data as a stratification of $Z$.

%Note $V$ is Henselian along its closed point by construction~\cite[Tag~0BSQ]{stacks-project} and so it follows that $\O_{X,V}^{\sh}$ is isomorphic to $\O_{X,t}^{\sh}$, hence $\piet(\beta)$ is also an isomorphism.
Using our stratification, we know that for any geometric point $x \in X$, the stalks of $R^1 j_* G_x$ at $s$ and $t$ are the same.
These stalks however are $H^1(T_s, G_x)$ and $H^1(T_t, G_x)$, respectively \cite[Chapter III, Theorem 1.15]{MilneEtaleCohomology}.
In other words, the functor of isomorphism classes of $G_x$-torsors over \autoref{eq.TubularMapOfHenselizations} produces isomorphisms.
But we have functorial bijections
\[ H^1(T_s, G_x) \cong \factor \Hom_{\cont} \big( \pi_1^{\et}(T_s), G_x \big). G_x. \]
and
\[ H^1(T_t, G_x) \cong \factor \Hom_{\cont} \big( \pi_1^{\et}(T_t), G_x \big). G_x., \]
where $G_x$ acts on $\Hom_{\cont}(-, G_x)$ by conjugation.
For more discussion see for instance \cite[Section 11.5]{GortzWedhornAlgebraicGeoemtryI}, \cite[Example 11.3]{MilneLecturesOnEtaleCohomology} and \cite[Chapter I, Remark 5.4 and Chapter III, Corollary 4.7, Remark 4.8]{MilneEtaleCohomology}.
%\todo{{\bf Karl:} Bhargav, is this right?  Do we have a better reference in this setting? \\
%{\bf Patrick:} I see this only if $G$ is abelian, unfortunately. Na\"\i vely, I think of an element of $H^1(X, G)$ as a \v Cech cocycle and then I construct a finite covering space by gluing. Then I want that covering space to be Galois with Galois group $G$, but this seems to be the case only when $G$ is abelian. Also \cite{MilneEtaleCohomology} seems to be stating the desired property only if $G$ is abelian (top of p.~121). \\
%{\bf Bhargav:} The identification of $H^1(X,G)$ (defined via Cech cocycles) with isomorphism classes of $G$-torsors on $X$ holds true for any topos $X$ and any sheaf $G$ of groups on $X$; one reference is \S 11.5 in G\"{o}rtz-Wedhorn. Also, if you further assume that $G$ is a constant sheaf associated to a finite group and $X$ is a connected scheme with a geometric base point $x$, then isomorphism classes of $G$-torsors are in bijection with conjugacy classes of maps $\pi_1(X,x) \to G$; the point is that any \'etale sheaf $F$ over $X$ which is a $G$-torsor is automatically representable by a finite \'etale cover of $X$ by flat descent for finite \'etale covers. Thus, all in all, conjugacy classes of maps $\pi_1(X,x) \to G$ identifies with $H^1(X,G)$ for $X$ connected. If you want to get the set of homomorphisms (not just modulo conjugation), then you need to work with isomorphism classes of $G$-torsors over $X$ equipped with a trivialization over $x$; however, I think we can just get by with conjugacy classes for the argument.}
In particular, the map
\begin{equation} \label{eq.alpha}
\factor \Hom \big( G_t \cong \piet(T_t), G_x \big). G_x. \longrightarrow \factor \Hom \big( G_s = \piet(T_s), G_x \big). G_x.
\end{equation}
induced by $\piet(\alpha)$ is a bijection for any geometric point $x \in X$.

Recall that we want to prove that $\piet(\alpha)$ is an isomorphism.
We apply~\autoref{eq.alpha} with $x = s$.
Considering the identity $\pi_1^{\et}(T_s) \to G_s$ as an element of the right-hand side, observe that there exists a homomorphism $\kappa\colon G_t \to G_s$, unique up to conjugacy, so that
\begin{equation} \label{eq.kappa}
G_s \xrightarrow{\piet(\alpha)} G_t \xrightarrow{\;\;\kappa\;\;} G_s
\end{equation}
is equal to a conjugate of the identity, that is, an inner automorphism.
In particular, $\kappa$ is surjective and $\pi_1^{\et}(\alpha)$ is injective.
We now apply $\Hom(-, G_t) / G_t$ to~\autoref{eq.kappa}. This gives us
\[ \factor \Hom(G_s, G_t). G_t. \longrightarrow \factor \Hom(G_t, G_t). G_t. \longrightarrow \factor \Hom(G_s, G_t). G_t.. \]
The second map, which is induced from $\piet(\alpha)$, is a bijection by~\autoref{eq.alpha} applied with $x = t$.
The composition is also a bijection, by construction.
Hence the first map, which is induced by $\kappa$, is bijective too.
By the same argument as above, this implies that there exists a homomorphism $\lambda\colon G_s \to G_t$ such that the composition $G_t \xrightarrow{\kappa} G_s \xrightarrow{\lambda} G_t$ is an inner automorphism.
In particular, $\kappa$ is injective and hence an isomorphism.
By~\autoref{eq.kappa}, it follows that $\piet(\alpha)$ is likewise an isomorphism, which was to be shown.
\end{proof}

\begin{remark} \label{rem.StratificationBehavesWellUnderPullback}
The crucial property of the stratification $\{ Z_i \}$ in \autoref{prop.GabberStratification} is preserved by finite quasi-\'etale covers:

Suppose that $\rho\colon X' \to X$ is a separated quasi-finite cover with branch locus contained in $Z$.
By Zariski's Main Theorem, we may assume that $\rho$ is finite since the statement is obvious for open inclusions.
Consider a specialization $s' \rightsquigarrow t'$ of geometric points of $\rho^{-1}(Z_i)$ mapping to a specialization of geometric points of $Z_i \subseteq X$, $s \rightsquigarrow t$, under $\rho$.
The diagram
\begin{equation} \label{fibre product}
\xymatrix{
\Spec \O_{X',s'}^{\sh} \ar[r] \ar[d]_-{\rho_s} & \Spec \O_{X',t'}^{\sh} \ar[d]^-{\rho_t} \\
\Spec \O_{X,s}^{\sh} \ar[r] & \Spec \O_{X,t}^{\sh}
}
\end{equation}
is a fibre product diagram up to taking a connected component (in the end it will follow that there can be at most one connected component from our assumptions).

Denote $T_{t'} := \Spec \O_{X', t'}^{\sh} \setminus \text{inverse image of $\rho^{-1}(Z)$}$ and likewise with $T_{s'}$.
Then removing the preimages of $Z$ from \eqref{fibre product} yields a fibre product diagram (up to taking a connected component, for now)
\[ \xymatrix{
T_{s'} \ar[rr]^-{\alpha'} \ar[d]_-{\rho_s} & & T_{t'} \ar[d]^-{\rho_t} \\
T_s \ar[rr]^-\alpha & & T_t.
} \]
The vertical maps are \'etale since the branch locus of $\rho$ is contained in $Z$, so the induced maps on $\pi_1^{\et}$ are injective.
Since $\piet(\alpha)$ is an isomorphism, we have that $\piet(T_{s'})$ and $\piet(T_{t'})$ define the same subgroups in $\piet(T_s) \cong \piet(T_t)$ (up to conjugation).
In particular, $\piet(\alpha')$ is also an isomorphism and so $\{ \rho^{-1}(Z_i) \}_{i \in I}$ stratifies $X'$ in the sense of \autoref{prop.GabberStratification}.
%
%Of course, strictly speaking $\{ \rho^{-1}(Z_i) \}$ is not a stratification of $X'$ since the $\rho^{-1}(Z_i)$ need not be irreducible.
%But we may stratify each $\rho^{-1}(Z_i)$ by declaring an arbitrary irreducible component $Z'_{i,1}$ of $\rho^{-1}(Z_i)$ to be a stratum, then picking an irreducible component of $\rho^{-1}(Z_i) \setminus Z'_{i,1}$ as the next stratum, and continuing in this fashion.
%This way we obtain a stratification $\{ Z'_{i,j} \}$ of $X'$ satisfying the conclusion of \autoref{prop.GabberStratification}.
%This stratification has the further property that the generic points of its strata are exactly the preimages of the generic points of the original stratification $\{ Z_i \}$ of $X$.
\end{remark}
%\todo{{\bf Patrick:} How about that?  {\bf Karl+Javier:} It looks ok to us, we made some minor changes.}

%\todo{{\bf Patrick:} I think more should be said here, e.g.~what is $T'_{t'}$ and where do we use that $\rho$ is quasi-\'etale?  {\bf Karl:} I added more explanation and notation but probably this should be said better. {\bf Patrick:} I'll try.  {\bf Karl:} Maybe another way to say this is as follows.  Take a finite universal cover over $T_t$.  Base change, and take a connected component, to get some Galois cover over $T_s$ (which is dominated by $T_s$'s finite universal cover a priori).  Since the induced map between the $\pi_1^{\et}$s is an isomorphism of finite groups, we must have that the pulled back cover is the universal cover, right (if the induced cover over $T_s$ is of smaller degree, there is no way the map of fundamental groups is going to be surjective)?  But now any given subcover under the universal cover of $t$ must base change to a cover (of the same degree) over $s$ by field degree nonsense, and hence we are done.  What do you think?\vskip 9pt
%Patrick has agreed to rewrite this more clearly.}

\begin{proposition} \label{prop.StrataAndCovers}
Suppose that $Z, U \subseteq X$ and let $\{ Z_i \}_{i \in I}$ be such a stratification as in \autoref{prop.GabberStratification} made up of connected $Z_i$.
Let $\rho\colon Y \to X$ be a finite surjective map whose branch locus is contained in $Z$ (in particular $\rho$ is quasi-\'etale), and let $W \subseteq X$ be the maximal Zariski-open subset over which $\rho$ is \'etale.
Then, for every $i \in I$, either $W \cap Z_i = \emptyset$ or $Z_i \subseteq W$.
Equivalently, either the branch locus of $\rho$ contains $Z_i$ or it is disjoint from $Z_i$.
\end{proposition}

\begin{proof}
We have the following criterion for geometric points of $X$ to belong to $W$.

\begin{claim} \label{clm.EtaleIffLocallyTrivial}
A geometric point $t \in X$ is contained in $W$ ($\rho$ is \etale{} over $t$) if and only if the pullback of $\rho\colon Y \to X$ to $T_t := \Spec \O_{X, t}^{\sh}\setminus Z^{\sh}$ is trivial (a finite disjoint union of copies of $T_t$).
\end{claim}

\begin{proof}[Proof of \autoref{clm.EtaleIffLocallyTrivial}]
Since the condition of being \'etale is Zariski-local, we see that $t \in W$ if and only if $\rho_{t}\colon Y_t = Y \times_{X} \Spec \O_{X,t} \to \Spec \O_{X,t}$ is \'etale (here we identify $t$ with its image at $X$ to avoid cumbersome phrasing and notation). However, by faithfully flat descent~\cite[Tag 02YJ, Lemma 34.20.29]{stacks-project} this morphism is \'etale if and only if its pullback to $\Spec \O_{X,t}^{\sh}$, say $\rho_t^{\sh}\colon Y_t^{\sh} \to \Spec \O_{X,t}^{\sh}$, is \'etale. But the target being the spectrum of a strictly henselian local ring, the latter condition implies that the finite \etale{} cover $\rho_t^{\sh}$ is trivial (\ie a finite disjoint union of copies of $\Spec \O_{X,t}^{\sh}$).
In particular, if $t \in W$ then the pullback $\rho^{-1}(T_t) \to T_t$ is trivial.

Conversely if $\rho^{-1}(T_t) \to T_t$ is trivial, then so is $\rho_t^{\sh}$ since\footnote{Note finite maps between normal schemes are determined outside a set of codimension 2 \cite{HartshorneGeneralizedDivisorsOnGorensteinSchemes}.} all the schemes are normal and the complement of $T_t \subseteq \Spec \O_{X,t}^{\sh}$ has codimension $\geq 2$.
As above this implies $t \in W$.
\end{proof}

%Let $t \in Z_i \cap W$, we want to show $t' \in W$ for any other geometric point $t' \in Z_i$. Let $s$ be the geometric generic point of $Z_i$. Notice that by assumption $T_t \to X$ factors trough the inclusion $W \cap U \to X$, in particular the pullback of $\rho$ to $T_t$ coincides with the pullback of its \etale{} restriction to $W \cap V$, say $\rho_{W \cap U}:Y_{W \cap W} \to W \cap U$.
Now let $t \in Z_i$ be a geometric point and let $s \in Z_i$ be a generic geometric point with a specialization $s \rightsquigarrow t$.
It is sufficient to show that $t \in W$ if and only if $s \in W$.
This follows by recalling that we have a canonical map $T_s \to T_t$ which induces an isomorphism on the level of fundamental groups, and moreover a commutative diagram
\[
\xymatrix@R=12pt{
T_t \ar[rd] \\
T_s \ar[u]\ar[r] & W \subseteq X. \\
}
\]
It follows that the images of $\piet(T_t)$ and of $\piet(T_s)$ coincide in $\piet(W)$ up to conjugation.
By~\cite[Chapter~I, Theorem~5.3]{MilneEtaleCohomology}, the finite \'etale cover $\rho^{-1}(W) \to W$ corresponds to a continuous $\piet(W)$-action on the finite set $Q := F(\rho^{-1}(W))$, where $F$ is the fibre functor.

Now assume $s \in W$, \ie~$\rho$ is \'etale over $s$.
By~\autoref{clm.EtaleIffLocallyTrivial},
% the induced action of $\piet(T_{s})$ on $Q$ is actually pulled back from $\piet(\Spec \O_{X,s}^\sh)$.
% But $\O_{X,s}^\sh$ is strictly henselian, so its $\piet$ is trivial, %{\color{blue} (right?)},
the induced action of $\piet(T_s)$ on $Q$ is trivial.
It follows that the induced action of $\piet(T_t)$ on $Q$ is also trivial.
By~\cite[Chapter I, Theorem 5.3]{MilneEtaleCohomology} again, $\rho^{-1}(T_t) \to T_t$, or more precisely $Y \times_X T_t \to T_t$, is trivial, \ie~the disjoint union of $\deg(\rho)$ copies of $T_t$.
Hence by \autoref{clm.EtaleIffLocallyTrivial} again, we get $t \in W$.
Conversely, $t \in W$ implies $s \in W$ by running the argument backwards.
\end{proof}

\section{Maximal quasi-\'etale covers}

As mentioned, our result is somewhat more general than the one in the introduction.

\begin{theorem}\textnormal{(\cf~\cite[Theorem 2.1]{GrebKebekusPeternellEtaleFundamental})} \label{thm.QuasiEtaleStabilize}
Suppose $X$ is an $F$-finite Noetherian integral strongly $F$-regular scheme.
Suppose that we have a commutative diagram of separated quasi-finite maps between normal $F$-finite Noetherian integral schemes
\[
\xymatrix{
    &              Y_1 \ar@{->>}[d]_{\eta_1}  & \ar[l]_{\gamma_1} Y_2 \ar@{->>}[d]_{\eta_2} & \ar[l]_{\gamma_2} Y_3 \ar@{->>}[d]_{\eta_3} & \ar[l]_{\gamma_3} \ldots\\
X   & \ar@{_{(}->}[l]^{j_0}   X_1                  & \ar@{_{(}->}[l]^{j_1}      X_2                 & \ar@{_{(}->}[l]^{j_2}      X_3                 & \ar@{_{(}->}[l]^{j_3} \ldots
}
\]
%\[ X = X_0 \xleftarrow{\;\;\gamma_1\;\;} X_1 \xleftarrow{\;\;\gamma_2\;\;} X_2 \xleftarrow{\;\;\gamma_3\;\;} \cdots \]
%is a sequence of finite quasi-\'etale covers by normal schemes, all of which are generically Galois over $X$.
such that the following conditions hold.
\begin{itemize}
\item[(i)] The maps $j_i$ are inclusions of open sets.
\item[(ii)] The maps $\eta_i$ are finite, quasi-\'etale and generically Galois.
\end{itemize}
Then all but finitely many of the $\gamma_j$ are \'etale.
\end{theorem}

To recover the statement in the introduction, simply set all the $X_i = X$.

\begin{proof}
By \cite[Theorem B]{deStefaniPolstraYao-Globalizing} (see also the semi-continuity of Hilbert-Kunz multiplicity and $F$-signature \cite{SmirnovUpperSemicontinuity,PolstraSemicontinuityFsignature,TuckerPolstraUniformApproach}) and the fact that $X$ is quasi-compact, we know there exists a uniform lower bound $\delta > 0$ on $s(\O_{X,x} =: R_x)$ for each scheme-theoretic point $x$ of $X$.
Let $R_x^{\sh}$ denote the strict henselization of the local ring $R_x$.
Since $R_x \subseteq R_x^{\sh}$ is flat with regular closed fiber, we know that $s(R_x^{\sh}) = s(R_x) \ge \delta$ by \cite[Theorem 5.6]{YaoObservationsAboutTheFSignature}.
Also $R_x^{\sh}$ is $F$-finite since quite generally, the strict henselization of a normal $F$-finite local ring is again $F$-finite\footnote{Indeed, if $R$ is an $F$-finite ring and $S$ is a direct limit of \etale{} maps (\ie an ind-\etale{} map), then $R^{1/p} \otimes_R S = S^{1/p}$ so that $S \subseteq S^{1/p}$ is finite.}.
In other words, we see that $s(\O_{X,x}^{\sh}) \geq \delta$ and $\O_{X,x}^{\sh}$ is $F$-finite for every geometric point $x$ of $X$.
By \cite[Theorem A]{CarvajalSchwedeTuckerEtaleFundFsignature}, we thus know that
\[
\big|\pi_1^{\et}(\Spec \O_{X,x}^{\sh} \setminus Z^{\sh})\big| \leq 1/\delta
\]
for each geometric point $x$ of $X$, where $Z$ is the singular locus of $X$.

Construct the stratification $\{ Z_i \}_{i \in I}$ as in \autoref{prop.GabberStratification}, where $Z = X_{\sing}$; assume the $Z_i$ are irreducible.  Observe that the $Z_i \cap X_k$ also stratify the open sets $X_k \subseteq X$ (although we may lose some pieces of the stratification).
Let $s_i$ be a geometric generic point of $Z_i$ and let $T_{s_i} := \Spec \O_{X, s_i}^{\sh} \setminus Z^{\sh}$ be its ``tubular neighborhood''. Pull back the whole sequence $X \leftarrow Y_1 \xleftarrow{\gamma_1} Y_2 \xleftarrow{\gamma_2} \cdots $ to $T_{s_i}$ to obtain
\[ T_{s_i} \xleftarrow{\;\;\gamma_{0, i}\;\;} T_{1, s_i} \xleftarrow{\;\;\gamma_{1, i}\;\;} T_{2, s_i} \xleftarrow{\;\;\gamma_{2, i}\;\;} \cdots. \]
Then all but finitely many of the $\gamma_{k, i}$ become trivial, as $\pi_1^{\et}(T_{s_i})$ is finite (of course, it is possible that $T_{k, s_i}$ becomes empty for $k \gg 0$ if $s_i \notin X_k$).
Note here is where we use the hypothesis that each $\eta_k : Y_k \to X_k \subseteq X$ is generically Galois, in particular Galois over the regular locus $U_k$ of $X_k$.

Now, given that $I$ is finite, one can pick $N \gg 0$ such that $\gamma_{n, s_i}$ is trivial for all $i \in I$ and all $n \geq N$. %(in other words that $\gamma_{n, s_i}$ is trivial for all $n \gg N$ and all $i$).
But by \autoref{rem.StratificationBehavesWellUnderPullback}, the inverse images of the generic points $s_i$ of the $Z_i$ in $Y_n$ are exactly the generic points $s_{n,j}$ of a stratification $\{ Z_{n,j} \}_{j \in J_n}$ of $Y_n$ satisfying the conclusion of \autoref{prop.GabberStratification}, with the $Z_{n,j}$ irreducible. Since
\[
\xymatrix{
\gamma_n\colon Y_{n+1} \ar@{->>}[r]^-{\text{finite}} & \eta_{n}^{-1}(X_{n+1}) \ar@{^{(}->}[r]^-{\text{open}} & Y_{n}
}
\]
is trivial after base changing with
\[ T_{s_{n,j}} := \Spec \O_{Y_n, s_{n,j}}^{\sh} \setminus \text{preimage of $\eta_n^{-1}(Z)$} \]
for $n \ge N$ and $j \in J_n$, by \autoref{clm.EtaleIffLocallyTrivial} we see that $\gamma_n$ is \'etale over the geometric generic point of every stratum $Z_{n,j}$ (of course, it might miss some completely).
Now applying \autoref{prop.StrataAndCovers} shows that $\gamma_n$ is finite \'etale over every point of $\eta_n^{-1}(X_{n+1})$, which completes the proof.
\end{proof}

The following corollary follows immediately.

\begin{corollary}\textnormal{(\cf~\cite[Theorem 1.5]{GrebKebekusPeternellEtaleFundamental})} \label{cor.MaximalQuasiEtale}
Suppose that $X$ is an $F$-finite Noetherian integral strongly $F$-regular scheme.
Then there exists a finite quasi-\'etale generically Galois cover $\rho\colon \tld X \to X$ with $\tld X$ normal and which satisfies the following property:

Every finite \'etale cover of the regular locus of $\tld X$ extends to a finite \'etale cover of $\tld X$.
Equivalently, the map $\piet(\tld X_{\reg}) \to \piet(\tld X)$ induced by the inclusion of the regular locus is an isomorphism.
\end{corollary}

We give two proofs of this result.
The first one uses \autoref{thm.QuasiEtaleStabilize}, while the second one is a direct proof which emphasizes the Galois correspondence.

\begin{proof}[First proof of \autoref{cor.MaximalQuasiEtale}]
Suppose not, then for every finite quasi-\'etale generically Galois cover $Y \to X$ there exists a further finite quasi-\'etale cover $Y' \to Y$, which we may also assume to be generically Galois over $X$, such that $Y' \to Y$ is not \'etale.
Repeating this process now with $Y'$ taking the role of $Y$, we obtain a sequence of covers contradicting the conclusion of \autoref{thm.QuasiEtaleStabilize}.

The proof of the statement about the map $\piet(\tld X_{\reg}) \to \piet(\tld X)$ is the same as~\cite[Step 2 of the proof of Thm.~1.5]{GrebKebekusPeternellEtaleFundamental}, and is thus omitted.
%Then, if $U := X_{\reg}$, there is a finite \'etale cover $U_1 \to U$ which does not extend across to $X$, this $U_1$ can be taken to be Galois over $U$ (by taking Galois closure if wanted). Let $X_1$ be the integral closure of $X$ in $U_1$ (or what is the same, in $K(U_1)$), thus we get a Cartesian square:
%\[\begindc{0}[5]
%\obj(-100, 100)[0]{$U_1$}
%\obj(50,0)[1]{$X$}
%\obj(50,100)[2]{$X_1$}
%\obj(-100,0)[3]{$U$}
%\mor{0}{2}{open}[1,6]
%\mor{0}{3}{$\text{Galois}$}[-1,0]
%\mor{2}{1}{$\gamma_1,\text{ finite}$}
%\mor{3}{1}{$\subseteq$}[1,6]
%\enddc
%\]
%Which exhibits $\gamma_1$ as a finite generically Galois cover of $X$, Galois over $U$ but not \'etale over $X$. Notice also that $U_1$ is contained in the regular locus of $X_1$ \cite[Chapter I, Proposition 3.17]{MilneEtaleCohomology}. We repeat this process for $X_1$ and its regular locus, to get $\gamma_2: X_2 \to X_1$ finite and Galois over the regular locus of $X_1$, so over $U_1$, but not \'etale over $X_1$, the composition $\gamma_2 \circ \gamma_1$ is then finite and Galois over $U$. This process can be repeated indefinitely contradicting \autoref{thm.QuasiEtaleStabilize}.
\end{proof}

\begin{proof} [Second proof of \autoref{cor.MaximalQuasiEtale}]
%By \cite[Theorem B]{deStefaniPolstraYao-Globalizing} (see also the semi-continuity of $F$-signature \cite{Smirnov, Polstra}) and the fact that $X$ is quasi-compact, we know there exists a uniform lower bound $\delta> 0$ on $s(\O_{X,x} =: R_x)$ for each scheme-theoretic point $x$ of $X$.  Let $R_x^{\sh}$ denote the strict henselization of the local ring $R_x$.  Since $R_x \subseteq R_x^{\sh}$ is flat with regular closed fiber \cite{}, we know that $s(R_x) = s(R_x^{\sh})$ by \cite{Yongwei}.  Since the strict henselization of a normal $F$-finite local ring is $F$-finite\footnote{Indeed, if $K$ is an $F$-finite field and $L$ is a separable extension, then $K^{1/p} \otimes_K L = (L)^{1/p}$ so that $L \subseteq L^{1/p}$ is finite.  Now let $K$ be the fraction field of $R_x$ and $L$ the fraction field of $R^{\sh}_x$, the result follows.}, we see that $s(\O_{X,x}^{\sh}) \geq \delta$ for every geometric point $x$ of $X$ as well.  By \cite{Carvajal-RojasSchwedeTucker}, we thus know that
%\[
%p \not\Big|\;\; \big|\pi_1^{\et}(\Spec \O_{X,x}^{\sh} \setminus Z^{\sh})\big| \leq 1/\delta
%\]
%for each geometric point $x$ of $X$.
Consider the stratification $\{ Z_i \}_{i \in I}$ as in \autoref{prop.GabberStratification} where each $Z_i$ is connected.  Set $Z = X_{\sing}$.
As before, for every $Z_i$ choose $s_i$ to be a geometric generic point.
As observed in the proof of \autoref{prop.StrataAndCovers}, if $t$ is a geometric point of $Z_i$ generizing to $s_i$ (as in \autoref{prop.GabberStratification}), then $\pi_1^{\et}(T_t)$ and $\pi_1^{\et}(T_s)$ have a common image in $\pi_1^{\et}(U)$, where $U = X \setminus Z$. %Further choose a $V$ witnessing this specialization as in  and remove the inverse image of $Z$ to obtain the maps between the $T$ as in \autoref{eq.TubularMapOfHenselizations}. All of these schemes map to $U$, and so $\pi_1^{\et}(T_t) \to \pi_1^{\et}(U)$ has a common image for each geometric $t \in Z_i$ specializing to $\eta_i$.
Let then $G_i \subseteq \pi_1^{\et}(U)$ denote this common image.
Note this common image is only unique up to conjugation.
In particular, since the index set $I$ is finite, there are only finitely many $G_i$ up to conjugation.

\begin{claim}
There is an open normal subgroup $H \leq \pi_1^{\et}(U)$ that intersects $G_i$ trivially for every $i$.%{\color{blue} Furthermore, we may choose that $p \not|\;\; [\pi_1^{\et}(U) : H]$.  (Or at least we will try...probably it is not true)}
\end{claim}

\begin{proof}[Proof of claim]
Note that $\pi_1^{\et}(U)$ is a profinite group
\[ \pi_1^{\et}(U) = \varprojlim Q_j \]
where the $Q_j$ are quotients of $\pi_1^{\et}(U)$.
Thus each $G_i$ maps injectively to $Q_j$ for large enough $j \gg 0$.
Choose a large enough $j$ that works for all $G_i$ and let $H = \ker(\pi_1^{\et}(U) \to Q_j)$.
In particular, $H$ is an open normal subgroup that intersects each $G_i$ trivially and hence by normality, it also intersects any conjugate trivially.
This proves the claim.
\end{proof}

Returning to the proof, by the Galois correspondence, there is a Galois cover $\varrho\colon \tld U \to U$ such that $\Aut_U(\tld U) = Q_j$ and $\pi_1^{\et}\bigl(\tld U\bigr) = H \subseteq \piet(U)$. %Let $\rho:\tld X \to X$ be the integral closure of $X$ in the fraction field of $\tld U$. %the unique integral closure of $\tld U \to U$, computed by Zariski's Main Theorem or the ``Main Theorem of Zariski''.
%\begin{comment}
%By Zariski's Main Theorem, we have a commutative square:
%\[\begindc{0}[5]
%\obj(-100, 100)[0]{$\tld U$}
%\obj(50,0)[1]{$X$}
%\obj(50,100)[2]{$\tld X$}
%\obj(-100,0)[3]{$U$}
%\mor{0}{2}{open}[1,6]
%\mor{0}{3}{$\varrho$}[-1,0]
%\mor{2}{1}{$\rho,\text{ finite}$}
%\mor{3}{1}{$\iota=\subseteq$}[1,6]
%\enddc
%\]
%where $\rho:\tld X \to X $ is actually the $\textbf{Spec}$ of the integral closure of $\sO_X$ in $(\iota \circ \varrho)_{*}\sO_{\tld U}$, therefore finite, and $\tld U$ maps as an open immersion into $\tld X$.
%\end{comment}
Let $\rho\colon \tld X \to X$ be the integral closure of $X$ in $K(\tld U)$, which is finite and its pullback to $U$ is exactly $\varrho$.
In particular, $\rho$ is quasi-\'etale and generically Galois.
In what remains, we prove that $\rho$ has the desired property. %From now onward we realize $\tld U$ as an actual open subscheme of $\tld X$ and think of $\varrho$ as the restriction of $\rho$ and denote them by the same letter $\rho$.

Let $V\to \tld U$ be a finite \'etale cover in $\textsf{F\'{E}t}\bigl(\tld U\bigr)$, we will extend it across $\tld X$. %\ie there exist $Y \to \tld X$ in $\textsf{F\'{E}t}\bigl(\tld X\bigr)$ which restricts to $V \to \tld U$.
Take $\sigma\colon Y \to \tld X$ to be the integral closure of $\tld X$ in $K(V)$. % the spectrum of the integral closure of $\O_{\tld X}$ in $K(V)$.
It suffices to prove that $Y \to \tld X$ is \'etale.

Let $W \subseteq \tld X$ be the complement of the branch locus of $\sigma$. Let $\tld Z$ denote the inverse image of $Z$ in $\tld X$ and note that $\tld U$ is the inverse image of $U$ in $\tld X$, so that $\tld U = \tld X \setminus \tld Z$.
Now $\tld U \subseteq W$ and we want to show that $W = \tld X$. %Then we have a criterion as in \autoref{clm.EtaleIffLocallyTrivial}, which somehow mixes it with \autoref{rem.StratificationBehavesWellUnderPullback}.
By \autoref{clm.EtaleIffLocallyTrivial}, a geometric point $x \in \tld X$ belongs to $W$ if and only if the pullback of $\sigma\colon Y \to \tld X$ to $\tld T_x := \Spec \O_{\tld X, x}^{\sh} \setminus \tld{Z}^{\sh}$ is trivial, where $\tld Z^{\sh}$ is the preimage of $\tld Z$.
%For this we establish a criterion for (geometric) points to be in $W$ and afterwards we argue how by construction this is satisfied for every point of $\tld X$.
%\begin{comment}
%\begin{claim}
%A geometric point $x \in \tld X$ belongs to $W$ if and only if the pullback of $\sigma:Y \to \tld X$ to $\Spec \O_{\tld X, x}^{\sh}\setminus \tld{Z}^{\sh}=:\tld T_{x}$ is trivial.
%\end{claim}
%\begin{proof}[Proof of claim] Since the condition of being \etale{} is Zariski-local, we see that $x \in W$ if and only if $\sigma_{x}:Y_x=Y\times_{\tld X} \Spec \O_{\tld{X},x} \to %\Spec \O_{\tld{X},x}$ is \etale{} (here we identify $x$ with its image at $\tld X$ to avoid cumbersome phrasing). However, by decent theory this morphism is \etale{} if and only if its pullback to $\Spec \O_{\tld{X},x}^{\sh}$, say $\sigma_x^{\sh}:Y_x^{\sh} \to \Spec \O_{\tld{X},x}^{\sh}$, is \etale{}. But being the target the spectrum of a strictly henselian local ring, the latter condition is equivalent to say that the finite cover $\sigma_x^{\sh}$ is trivial (\ie finite disjoint union of copies of $\Spec \O_{\tld{X},x}^{\sh}$). However, this triviality can be checked by restricting to an open subset% whose complement has codimension at least two
%\footnote{Under suitable hypothesis, say $U \subset X$ a connected dense open subscheme of a connected normal locally excellent scheme $X$, the restriction functor $\text{F\'{E}t}(X) \to \text{F\'{E}t}(U)$ is fully faithful \cite[OBQF]{stacks-project}.}, for example to the tube $\tld T_{x}:=\Spec \O_{\tld X,x}^{\sh} \setminus \tld{Z}^{\sh}$.
%This proves the claim.
%\end{proof}
%\end{comment}

%Because $\tld Z$ contains the complement of $W$, so
For any geometric point $x \in \tld X$, $\tld T_x \to \tld X$ factors through $\tld U$.
Thus the pullback of $\sigma\colon Y \to \tld X$ to $\tld T_x$ coincides with the pullback of $V \to \tld U$ to $\tld T_x$.
We see that $x \in W$ if and only if the pullback of $V \to \tld U$ to $\tld T_x$ is trivial.
Hence it suffices to show that for all geometric points $x \in \tld X$ and all $V / \tld U \in \textsf{F\'{E}t}(\tld U)$, the pullback of $V \to \tld U$ to $\tld T_x$ is trivial.
Equivalently, we want the induced homomorphism of fundamental groups $\piet \bigl( \tld T_x \bigr) \to \piet \bigl( \tld U \bigr)$ to be zero~\cite[5.2.3]{MurreLecturesFundamentalGroups}. %, this for all $x \in X$.
This is argued below.

Let $t \in X$ be the image of the geometric point $x \in \tld X$. We have a commutative diagram
\[
\xymatrix{
\tld T_x \ar[r] \ar[d] & \tld U \ar[d] \\
T_t \ar[r] & U
}
\]
%\[\begindc{0}[5]
%\obj(-100, 100)[0]{$\tld T_x$}
%\obj(0,0)[1]{$U$}
%\obj(0,100)[2]{$\tld U$}
%\obj(-100,0)[3]{$T_t$}
%\mor{0}{2}{}
%\mor{0}{3}{}
%\mor{2}{1}{}
%\mor{3}{1}{}
%\enddc
%\]
with finite \'etale vertical morphisms.
By applying $\piet(-)$ we obtain the commutative square
\[
\xymatrix{
\pi_1^{\et}\bigl(\tld T_x\bigr) \ar@{^{(}->}[d] \ar[r] & H \ar@{^{(}->}[d] \\
\pi_1^{\et}(T_t) \ar[r] & \pi_1^{\et}(U)
}
\]
%\]
%\[\begindc{0}[5]
%\obj(-100, 100)[0]{$\pi_1^{\et}\bigl(\tld T_x\bigr)$}
%\obj(50,0)[1]{$\pi_1^{\et}(U).$}
%\obj(50,100)[2]{$H$}
%\obj(-100,0)[3]{$\pi_1^{\et}(T_t)$}
%\mor{0}{2}{}
%\mor{0}{3}{}[1,6]
%\mor{2}{1}{}[1,6]
%\mor{3}{1}{}
%\enddc
%\]
However, by construction $H$ meets the image of $\pi_1^{\et}(T_t)$ in $\pi_1^{\et}(U)$ trivially, which forces the top map in this square to be zero, as required.
%By \cite{Milne's Ch I prop 3.8} there is a maximal Zariski-open of $Y$ on which this map is \etale{}, since it is finite (in particular closed) the image of its complement is closed, whose open complement $W \subseteq \tld X$ is the maximal Zariski-open of $\tld X$ on which $\sigma: \sigma^{-1} W \to W$ is \etale{}. Notice $W$ contains $\tld U$, we want to show $W=\tld X$, for this purpose define $Z'$ to be the complement of $\tld U$ in $\tld X$. By decent theory $x \in W=W_{\sigma}$ (a geometric point) if and only if $o \in W=W_{\sigma'}$ for some (then all) \etale{} neighborhood $(O,o)$ ($o$ geometric point) of $x$, where $\sigma'$ is the induced morphism $Y_O \to O$.
\end{proof}

%\todo{{\bf Patrick:} I have made some minor edits to the above proof (some $U$s should be $\tld U$s and so on).}

If $X$ is a strongly $F$-regular proper variety over an $F$-finite field $k$, and $U = X_{\reg}$, then it follows from \cite[Theorem 7.6]{SchwedeTuckerTestIdealFiniteMaps} that any \etale{} cover $U$ is cohomologically tame and hence tame in all the senses of \cite{KerzSchmidtOnDifferentNotionsOfTameness}.
%In particular, $\pi_1^{\text{t}}(U) = \pi_1^{\et}(U)$ {\color{blue} {\bf Karl:} right?}.
From this we immediately obtain the following corollary.

\begin{corollary} \label{cor.MaximalTameCover}
If $X$ is a strongly $F$-regular proper variety over an $F$-finite field $k$ with $U = X_{\reg}$, then the map $\tld X \to X$ from \autoref{cor.MaximalQuasiEtale} restricts to a tame \etale{} cover of $U$.
\end{corollary}

We also obtain a local version of \autoref{cor.MaximalQuasiEtale}.

\begin{corollary} \textnormal{(cf.~\cite[Theorem 1.9]{GrebKebekusPeternellEtaleFundamental})} \label{cor.LocalMaximalQuasiEtale}
Suppose that $X$ is an $F$-finite Noetherian integral strongly $F$-regular scheme and that $x \in X$ is a point.
%Let $U = X_{\reg}, Z = X_{\sing} \subseteq X$ be as in \autoref{cor.MaximalQuasiEtale}.
Then there exists a Zariski-open neighborhood $x \in X^\circ \subseteq X$ and a finite quasi-\'etale generically Galois cover $\rho\colon \tld X^\circ \to X^\circ$ with $\tld X^\circ$ normal and which satisfies the following property:

For every further Zariski-open neighborhood $x \in W \subseteq X^\circ$, with $\tld W = \rho^{-1}(W)$, any further finite quasi-\'etale cover $\overline{W} \to \tld W$ is \'etale.
Equivalently, $\piet(\tld W_{\reg}) \to \piet(\tld W)$ is an isomorphism.
\end{corollary}

\begin{proof}
Suppose not, \ie~assume that for every $x \in X^\circ \subseteq X$ and every finite quasi-\'etale generically Galois $\tld{X^\circ} \xrightarrow\rho X^\circ$, there is $x \in W \subseteq X^\circ$ and a finite quasi-\'etale $\overline W \to \rho^{-1}(W)$ that is not \'etale.
By taking Galois closure, we may assume that $\overline W \to W$ is Galois.

Apply this assumption with $X^\circ := X$ and $\rho := \id_X$.
We obtain a map $\overline W \to W$, which we denote $\eta_1\colon Y_1 \to X_1$.
Applying the assumption again, this time to $\eta_1$, we get a map $\eta_2\colon Y_2 \to X_2$ together with maps $\gamma_1\colon Y_2 \to Y_1$ and $X_2 \hookrightarrow X_1$.
Inductively, we construct a diagram as in \autoref{thm.QuasiEtaleStabilize}, but where none of the $\gamma_i$ are \'etale.
This contradicts \autoref{thm.QuasiEtaleStabilize}.

%Suppose not, then choose $\eta_1\colon Y_1 \to X$ to be a cover satisfying \autoref{cor.MaximalQuasiEtale} and set $X_1 = X$.  By our hypothesis (which we intend to contradict), there exists a Zariski-open neighborhood $X_2$ of $x$ and a quasi-\etale{} cover $V_2 \to X_2$ (which we may assume is generically Galois) which is not \etale.  Let $Y_1' = \eta_1^{-1}(X_2)$ and choose a finite quasi-\etale{} cover $Y_2 \to X_2$ dominating $Y_1'$ and $V_2$ satisfying \autoref{cor.MaximalQuasiEtale}.  Continuing in this way contradicts \autoref{thm.QuasiEtaleStabilize}.
\end{proof}

\begin{remark}
Note that one cannot always take $X^{\circ} = X$ in the statement of \autoref{cor.LocalMaximalQuasiEtale}, even in characteristic zero.
Indeed let $X$ be the \emph{projective} quadric cone $V(x^2 - yz) \subseteq \bP^3_{\bC}$.
Then $X_{\reg}$ is simply connected (it is an $\mathbb A^1$-bundle over $\bP^1_{\bC}$), so every finite quasi-\'etale cover of $X$ is trivial.
On the other hand, the \emph{affine} quadric cone $U \subseteq X$ does have a finite quasi-\'etale cover that is not \'etale, corresponding to $k[x^2, xy, y^2] \subseteq k[x,y]$.
%: a ruling $D \subseteq U$ is a Weil divisor that is not Cartier, but $2 \cdot D$ is linearly equivalent to zero.

In characteristic 2, $\pi_1^{\et}$ of the punctured quadric cone singularity is trivial and hence this computation does not work (the local cover above is in fact inseparable).
However, one can obtain the same conclusion over an algebraically closed field of characteristic $p > 2$.
Let $X' \to X$ be a finite quasi-\'etale generically Galois cover of $X$ that is not \'etale.
Let $O \in X$ be the cone point.
It follows that
\[ X' \times_X \Spec \O_{X, O}^{\sh} \]
is a disjoint union of copies of $V \to \O_{X,O}^{\sh}$ where $V$ is the regular 2-to-1 cover of $\Spec \O_{X, O}^{\sh}$ (corresponding to $k[x^2, xy, y^2] \subseteq k[x,y]$).

Let $L$ be a ruling of the cone and consider $X' \times_X L \to L$.  This map is \etale{} except over the cone point $O \in L \subseteq X$, and over that point $X' \times_X L$ is not even reduced.  However, if one takes $L' = (X' \times_X L)_{\red}$ then the computation $k\llbracket x^2, xy, y^2 \rrbracket/\langle x^2, xy \rangle \subseteq k\llbracket x,y \rrbracket /\langle x^2, xy \rangle$ shows that $L'$ is normal and furthermore that $L' \to L$ is ramified of order 2 over $O$ on each connected component.  In particular, $L' \to L$ is ramified over a single point and that ramification is tame (since $p > 2$), but $L \cong \bP^1$ and so this is a contradiction.
Hence any quasi-\etale{} cover of $X$ is in fact \etale{} and so we cannot take $X^{\circ} = X$ in \autoref{cor.LocalMaximalQuasiEtale} just as in characteristic zero.
\end{remark}

\begin{remark}
We believe one can prove \autoref{cor.LocalMaximalQuasiEtale} using a strategy similar to \autoref{cor.MaximalQuasiEtale}.  In particular, we can first assume that $x$ is in the closure of any stratum $Z_i$ (if not, shrink $X$ to remove those strata).  Then use
\[
\pi_1^{\et}\big((\Spec \O_{X,x}) \setminus \text{(inverse image of $Z$)}\big)
\]
as the replacement for $\pi_1^{\et}(U)$ in the proof of \autoref{cor.LocalMaximalQuasiEtale}.  The $H$ we obtain produces a cover that is quasi-\'etale over a neighborhood $X^\circ$ of $x$ and which satisfies the desired property.
%\todo{{\bf Karl:} Does this look ok to everyone? {\bf Patrick:} Probably it is, but we should discuss it tomorrow.}
\end{remark}

Just as in \cite[Theorem 1.10]{GrebKebekusPeternellEtaleFundamental}, we obtain a result on simultaneous index-one covers.

\begin{corollary} \label{cor.SimultaneousIndex1Cover}
Suppose that $x \in X$ and $\tld X^\circ \twoheadrightarrow X^\circ \subseteq X$ is as in \autoref{cor.LocalMaximalQuasiEtale}.
Suppose further that $X^\circ$ is chosen to be affine (or quasi-projective over an affine scheme).
\begin{itemize}
\item[(a)]  If $\tld D$ is any $\bZ_{(p)}$-Cartier\footnote{This just means it is $\bQ$-Cartier with index not divisible by $p$.} divisor on $\tld X^\circ$, then $\tld D$ is Cartier.
\item[(b)]  There exists an integer $N > 0$ so that if $D$ is any $\bZ_{(p)}$-Cartier divisor on $X^\circ$, then $N \cdot D$ is Cartier.
\end{itemize}
\end{corollary}

\begin{proof}
The proof is exactly the same as the proof of \cite[Theorem 1.10]{GrebKebekusPeternellEtaleFundamental}, and so we omit it.
Of course, the main idea is to take a cyclic cover.
The hypothesis that $X^\circ$ is affine (or more generally quasi-projective over an affine) is assumed so that for any finite collection of points $S \subseteq X^\circ$ and for any line bundle $\sL$ on $X$, that there exists an open neighborhood of $S$ which trivializes $\sL$.
\end{proof}

Just as in \cite{GrebKebekusPeternellEtaleFundamental}, \autoref{cor.SimultaneousIndex1Cover} implies that for any $F$-finite Noetherian integral strongly $F$-regular scheme, there exists an integer $N > 0$ such that if $D$ is $\bZ_{(p)}$-Cartier, then $N \cdot D$ is Cartier.
The point is that by quasi-compactness, $X$ can be covered by finitely many open $X^\circ$ satisfying \autoref{cor.SimultaneousIndex1Cover}.

%\begin{remark}
%We believe it is possible to also recover the main result of \cite{GrebKebekusPeternellEtaleFundamental}, in characteristic zero, via the method of this paper.  Indeed, proceed as follows.
%\todo{{\bf Karl:} Bhargav, did you want to write something along these lines?}
%\end{remark}

\bibliographystyle{skalpha}
\bibliography{MainBib}

\end{document}
